\begin{Answer}{5.1}
Đặt $a=BC, b=CA, c=AB$ và $s = \frac{{a + b + c}}{2}$. Các tam giác ABH và ACH đồng dạng với tam giác ABC với tỉ số tương ứng $a/c$ và $b/c$\\
Áp dụng công thức diện tích tam giác bằng bán kính đường tròn nội tiếp nhân với nửa chu vi, suy ra bán kính cần tìm là $\frac{{ab}}{{a + b + c}};\frac{a}{c}\frac{{ab}}{{a + b + c}};\frac{b}{c}\frac{{ab}}{{a + b + c}}$ và tổng của chúng là $\frac{{ab}}{c} = AH$
\end{Answer}
\begin{Answer}{5.2}
Lưu ý rằng $a_{n + 1}^2  + b_{n + 1}^2  = \left( {\alpha ^2  + \beta ^2 } \right)\left( {a_n^{^2 }  + b_n^{^2 } } \right)$, trừ $\alpha  = \beta  = 0$. Chúng ta cần $\alpha ^2  + \beta ^2  = 1$. Vì vậy có thể đặt $\alpha  = \cos \theta ,\beta  = \sin \theta $, từ đó bằng phương pháp quy nạp ta chỉ ra $\alpha  = \cos n \theta ,\beta  = \sin n \theta $. Từ đó có 1998 bộ số: $\left( {0;0} \right)$ và $\left( {\cos \theta ;\sin \theta } \right)$ với $\theta  = \pi \frac{k}{{3998}},k = 1,3,...,3997$
\end{Answer}
\begin{Answer}{5.3}
Chúng ta thừa nhận rằng số tiền lớn nhất kiếm được bởi không bao giờ cho phép một cầu thủ chuyển đến đội nhỏ hơn. Chúng ta cũng có thể giữ kỷ lục đó bằng một cách khác. Một đội bóng có $x$ cầu thủ thì được ghi là $-x$ trước khi giao dịch một cầu thủ hoặc $x$ trước khi nhận một cầu thủ và số tiền mà liên đoàn kiếm được bằng tổng của các số đó. Bây giờ ta xem xét các số được ghi bởi một đội  mà kết thúc có nhiều hơn 20 cầu thủ. Nếu số lượng cầu thủ tối đa của đội đó trong suốt quá trình chuyển nhượng là $k>n$  thì các số $k-1$ và $-k$ xuất hiện liên tiếp và bỏ đi 2 số đó thì tổng sẽ tăng lên. Vì thế tổng của các số trong đội bóng đó ít nhất là $20+21+...+(n-1)$.\\
Tương tự như vậy, tổng của các số trong đội bóng kết thúc có $n<20$ cầu thủ ít nhất là $-20-19-...-(n+1)$. Vì những con số này chính xác là những con số có được bởi việc luân chuyển cầu thủ từ đội kết thúc ít hơn 20 cầu thủ sang đội kết thúc có nhiều hơn 20 cầu thủ. Sự sắp xếp đó dẫn đến số tiền kiếm được là lớn nhất. Trong trường hợp đó, tổng là:
$$ (20+20+21+20+21+20+21+22) -2(20+19+18+17)=17$$
\end{Answer}
\begin{Answer}{5.4}
Nếu gọi $2\alpha ,2\beta ,2\gamma ,2\delta $ là các cung chắn bởi các cạnh $a,b,c,d$ tương ứng thì:
$$ AC = 2\sin \left( {\alpha  + \beta } \right) = \frac{a}{2}\sqrt {1 - \frac{{b^2 }}{4}}  + \frac{b}{2}\sqrt {1 - \frac{{a^2 }}{4}} $$
Tương tự với $CD$.\\
Tổng quát với $R$ là bán kính đường tròn ngoại tiếp ngũ giác, thế thì $AC^2+BD^2=1$ mà
$$  AC = a\sqrt {R^2  - b^2 }  + b\sqrt {R^2  - a^2 } $$
Khi đó, dẫn đến biểu thức chứa $R^2$ dưới dấu căn và ta giải phương trình đối với $R$ theo các số $a,b,c,d$.
\end{Answer}
\begin{Answer}{5.5}
 Gọi $n$ là cấp của $5/3 \mod p$, và đặt $x_i=3^{n-1-i}5^{i-1}, y_i=43^{n-1}5^{i-1}$ thì mọi đồng dư trên là tương đương trừ hệ thức cuối cùng. Hệ thức đó có dạng $5^{2n}  \equiv 3^{2n} \left( {\bmod p} \right)$ (đúng).Vậy ta có điều phải chứng minh.
\end{Answer}
\begin{Answer}{5.6}
 Đặt $S=x_1+x_2+...+x_n$\\
Sử dụng BĐT Chebyshev's cho hai dãy $\frac{{x_i }}{{S - x_i }}$ và $x_i^4 $ (cả hai dãy đều là dãy tăng). Ta có:
$$ \sum {\frac{{x_i^5 }}{{S - x_i }} \ge \left( {\frac{{x_1 }}{{S - x_1 }} + \frac{{x_2 }}{{S - x_2 }} + ... + \frac{{x_n }}{{S - x_n }}} \right)} \left( {\frac{{x_1^4  + x_2^4  + ... + x_n^4 }}{n}} \right) $$
Áp dụng bất đẳng thức hàm lồi ta có:
$$\frac{{x_1 }}{{S - x_1 }} + \frac{{x_2 }}{{S - x_2 }} + ... + \frac{{x_n }}{{S - x_n }} \ge \frac{1}{{n - 1}}  $$
Áp dụng bất đẳng thức giá trị trung bình ta có:
$$ \left( {\sum {\frac{{x_i^4 }}{n}} } \right)^{{\raise0.7ex\hbox{$1$} \!\mathord{\left/
 {\vphantom {1 2}}\right.\kern-\nulldelimiterspace}
\!\lower0.7ex\hbox{$2$}}}  \ge \sum {\frac{{x_i^2 }}{n}}  = \frac{1}{n} $$
Ta có kết luận:
$$ \sum {\frac{{x_i^5 }}{{S - x_i }}}  \ge n\frac{1}{{n - 1}}.\frac{1}{{n^2 }} = \frac{1}{{n\left( {n - 1} \right)}} $$
Đẳng thức xảy ra khi $x_1  = x_2  = ... = x_n  = \frac{1}{{\sqrt n }}$.
\end{Answer}
